\scanpage{464}\mspage{376}
% archive sheet number: 463
\uline{Démonstration} Soit $\pi'_0$ \deleted{\unsure{un cardinal}} \inserted{$\unsure{\geq|>}\pi_0$ un} cardinal utile pour $M$, $\pi$ un cardinal ayant les propriétés suivantes
\[
\begin{aligned}
\text{a)}\quad &\pi\geq\operatorname{Sup}(\pi'_0,\operatorname{card\,red}M_{\pi'_0}),\quad\pi>\pi_0\\
\text{b)}\quad &\pi^{\pi'_0}=\pi\\
\text{c)}\quad &\underline I\in\deleted{\unsure{\Phi\cup\Psi}}\inserted{\Phi_0\cup\Psi}\Longrightarrow
\operatorname{card\,red}\underline I\leq\pi .
\end{aligned}
\]
\deleted{\unsure{$\mathrm{avec}\ \Phi$} \illegible{2} \unsure{$I$}} \deleted{On voit} Soit $C=M_\pi$, d'où $M\underset{\unsure{-|\leftarrow}}{\approx}\mathrm{Ind}_\pi(C)$. Dans tous les cas envisagés, $\Phi$ contient l'ens des cat. \inserted{\unsure{ess.|ens.|équ.} petites} qui sont gdes devant $\pi_0$, \unsure{a|à} fortiori celles qui sont gdes devant $\pi$, donc la catégorie $\underline{\mathrm{Hom}}_\Phi^\Psi$ \deleted{\unsure{est}} \inserted{s'interprète comme} une sous-catégorie \unsure{strictmt|strictement} pleine de
\[
\underline{\mathrm{Hom}}(C,N).
\]
Si d'autre part $\Phi_\pi$ désigne l'ens des $I\in\Phi$ telles que $\operatorname{card} I\leq\unsure{\pi|1}$, on trouve
\[
\underline{\mathrm{Hom}}_\Phi^\Psi(M,N)\underset{\unsure{-|\to}}{\approx}
\underline{\mathrm{Hom}}_{\Phi_\pi}^\Psi(C,N),
\]
et on applique 4.14.2.
